\documentclass[a4paper,12pt]{article}
\setcounter{secnumdepth}{5}
\setcounter{tocdepth}{3}
\input{/usr/share/latex-toolkit/template.tex}
\begin{document}
\title{Measure Theory}
\author{沈威宇}
\date{\temtoday}
\titletocdoc
\section{Measure Theory (測度論)}
\ssc{\text{\textsigma}-algebra or \text{\textsigma}-field}
\sssc{\subsection{\text{\textsigma}-algebra or \text{\textsigma}-field}
Let $X$ be some set, and let $\calP(X)$ represent its power set. A set $\Sigma \subseteq\calP(X)$ is called a \text{\textsigma}-algebra or \text{\textsigma}-field if and only if it satisfies the following three properties:
\begin{enumerate}
\item $X\in\Sigma$.
\item Closed under complementation:
\[A\in\Sigma\implies X\setminus A\in\Sigma.\]
\item Closed under countable unions: For all indexed family $(A_i)_{i\in\bbN}\subseteq\Sigma$,
\[\bigcup_{i\in\bbN}A_i\in\Sigma.\]
\epr
These imply:
\[\varnothing\inSigma\]
\bpr
\[X\in\Sigma\implies X\setminus X=\varnothing\in\Sigma.\]
\epr
These imply closed under finite unions: For all indexed family $(A_i)_{i\in\bbN\cap[1,n]}\subseteq\Sigma$,
\[\bigcup_{i=1}^nA_i\inSigma.\]
\bpr
Let $(B_i)_{i\in\bbN}$ be defined as $B_i=A_i$ for $i\in[1,n]$ and $B_i=\varnothing$ for $i\in(n,\infty)$ and apply closed under countable unions.
\epr
\end{enumerate}
\sssc{Measurable space (可測空間)}
A tuple $(X,\Sigma)$ of a set $X$ and a \text{\textsigma}-algebra $\Sigma$ on $X$ is called a measurable space.
\sssc{Intersection is \text{\textsigma}-algebra}
Let $X$ be some set, and let $P(X)$ represent its power set. If $\calF$ is a set of \text{\textsigma}-algebras of $X$, then $\bigcap_{\Sigma\in\calF}\Sigma$ is a \text{\textsigma}-algebra of $X$.
\bpr
If $A\in\bigcap_{\Sigma\in\calF}\Sigma$, then $\forall\Sigma\in\calF\colon A\in\Sigma$, then $\forall\Sigma\in\calF\colon X\setminus A\in\Sigma$, then $X\setminus A\in\bigcap_{\Sigma\in\calF}\Sigma$.

If $\calS\subseteq\bigcap_{\Sigma\in\calF}\Sigma$ and $\abs{\mathcal{S}}\in\bbN\cup\{\aleph_0\}$, then $\forall\Sigma\in\calF\colon\calS\subseteq\Sigma$, then $\forall\Sigma\in\calF\colon\bigcup_{A\in\mathcal{S}}A\in\Sigma$, then $\bigcup_{A\in\mathcal{S}}A\in\bigcap_{\Sigma\in\calF}\Sigma$.
\epr
\sssc{\text{\textsigma}-algebra generated by collection of subsets}
For a set $X$ and a collection $C\subseteq \calP(X)$, the \text{\textsigma}-algebra generated by $C$, denoted as $\sigma(C)$ or $\sigma\{C\}$, is the smallest \text{\textsigma}-algebra that is a superset of $C$.
\sssc{Borel \text{\textsigma}-algebra}
Let $X$ be a topological space, the Borel \text{\textsigma}-algebra of it, $\calB(X)$, is the \text{\textsigma}-algebra generated by the set of all Borel sets of $X$.
\sssc{Borel measurable space}
A tuple $(X,\calB)$ of a set $X$ and a Borel \text{\textsigma}-algebra $\calB$ on $X$ is called a Borel measurable space.
\sssc{Intervals generating Borel \text{\textsigma}-algebra of the real numbers theorem}
The Borel \text{\textsigma}-algebra of the real numbers can be generated by the collection of all intervals of the form $(-\infty,x]$ for $x\in\bbR$.
\bpr
For every $a<b$, the interval $(a,b]$ can be expressed as
\[(a,b]=(-\infty,b]\cup(-\infty,a]^c.\]
For every $a<b$, the interval $(a,b)$ can be expressed as
\[(a,b)=\bigcup_{n=1}^{\infty}(a,b-\flatfrac{1}{n}].\]
\epr
\sssc{\textpi-system or pi-system}
A \textpi-system or pi-system on a set $\Omega$ is a set $P$ of subsets of $\Omega$ such that
\bit
\item $P\ne\varnothing$, and
\item $A,B\in P\implies A\cap B\in P$.
\eit
\sssc{Generated \textpi-system}
For any non-empty set $\Sigma$ of subsets of a set $\Omega$, there exists a \textpi-system $P$, called \textpi-system generated by $\Sigma$, that is the unique smallest \textpi-system of $\Omega$ containing every element of $\Sigma$. It is equal to the intersection of all \textpi-systems of $\Omega$ containing every element of $\Sigma$.
\bpr
PLACEHOLDER
\epr
\sssc{\textlambda-system, lambda-system, Dynkin system, or d-system}
A \textlambda-system, lambda-system, Dynkin system, or d-system on a set $\Omega$ is a set $D$ of subsets of $\Omega$ such that:
\bit
\item $\Omega\in D$,
\item Closed under complements of subsets in supersets: $A,B\in D\land A\subseteq B\implies B\setminus A\in D$, and
\item Closed under countable increasing union: if $A_1\subseteq A_2\subseteq\ldots$ is an increasing sequence of sets in $D$, then $\bigcup_{i=1}^{\infty}A_n\in D$.
\eit
\sssc{\textpi-system, \textlambda-system, and \textsigma-algebra}
\ben
\item A set is both a \textpi-system and a \textlambda-system iff it is a \textsigma-algebra.
\item A \textpi-system is not necessarily a \textlambda-system.
\item A \textpi-system is not necessarily a \textsigma-algebra.
\item A \textlambda-system is not necessarily a \textpi-system.
\item A \textlambda-system is not necessarily a \textsigma-algebra.
\een
\bpr
PLACEHOLDER
\epr
\sssc{(Sierpiński–Dynkin's) \textpi-\textlambda theorem}
If $P$ is a \textpi-system and $D$ is a \textlambda-system with $P\subseteq D$, then the \textsigma-algebra generated by $P$ is a subset of $D$.
\bpr
PLACEHOLDER
\epr
\ssc{Measure}
\sssc{Measure (測度)}
Let $(X,\Sigma)$ be a measurable space. A function $\mu\colon\Sigma\to[-\infty,+\infty]$ is called a (positive) measure on $\Sigma$ iff it satisfies
\bit
\item Countable additivity, also called \text{\textsigma}-additivity: For all indexed family $(E_k)_{k\in\bbN}$ of pairwise disjoint sets in $\Sigma$,
\[\mu\left(\bigcup_{k\in\bbN}E_k\right)=\sum_{k=1}^\infty\mu(E_k).\]
\item Nonnegativity:
\[\forall E\in\Sigma\colon\mu(E)\geq 0,\]
\eit
Countable additivity implies $\mu(\varnothing)=0$.
\bpr
By applying for $(\varnothing)_{k\in\bbN}$,
\[\mu(\varnothing)=\sum_{k=1}^{\infty}\mu(\varnothing)\]
\epr
Countable additivity implies finite additivity: For all indexed family $(E_k)_{k\in\bbN\cap[1,n]}$ of pairwise disjoint sets in $\Sigma$,
\[\mu\left(\bigcup_{k=1}^nE_k\right)=\sum_{k=1}^n\mu(E_k).\]
\bpr
Let $(F_k)_{k\in\bbN}$ be defined as $F_k=E_k$ for $k\in[1,n]$ and $F_k=\varnothing$ for $k\in(n,\infty)$ and countable additivity.
\epr
If $\mu(E)\geq 0$ is replaced with $\mu(E)\in\bbR$, it's called a signed measure.

If $\mu(E)\geq 0$ is replaced with $\mu(E)\in\bbC$, it's called a complex-valued measure.
\sssc{Measure space (測度空間)}
A tuple $(X,\Sigma,\mu)$ of a set $X$, a \text{\textsigma}-algebra $\Sigma$ on $X$, and a measure $\mu$ (sometimes signed) on $\Sigma$ is called a measurable space.
\sssc{Pre-measure}
Let $X$ be a set and $\calR$ be ring of subsets on $X$. A function $\mu\colon\calR\to[-\infty,+\infty]$ is called a signed pre-measure on $\calR$ iff it satisfies
\bit
\item $\mu(\varnothing)=0$ and
\item countable additivity, also called \text{\textsigma}-additivity, that, for all finite or countable infinite indexed family $(E_k)_{k\in K}$ of pairwise disjoint sets in $\calR$,
\[\mu\left(\bigcup_{k\in K}E_k\right)=\sum_{k=1}^\infty\mu(E_k).\]
\eit
If it also satisfies nonnegativity:
\[\forall E\in\calR\colon\mu(E)\geq 0,\]
it is called a pre-measure or positive pre-measure.
\sssc{Outer or exterior measure}
Let $X$ be some set, and let $\calP(X)$ represent its power set. A function $\mu\colon\calP(X)\to[-\infty,+\infty]$ is called a signed outer or exterior measure on $\calR$ iff it satisfies
\bit
\item $\mu(\varnothing)=0$ and
\item countable subadditivity, that, for all finite or countable infinite indexed family $(E_k)_{k\in K}$ of pairwise disjoint sets in $\calP(X)$,
\[A\subseteq\bigcup_{k\in K}E_k\right)\implies\mu(A)\leq\sum_{k=1}^\infty\mu(E_k).\]
\eit
If it also satisfies nonnegativity:
\[\forall E\in\calR\colon\mu_0(E)\geq 0,\]
it is called a outer or exterior measure or outer or exterior measure.
\sssc{Product measure space}
Let $(X,\mathcal{A},\mu),(Y,\mathcal{B},\nu)$ be two measure space with possibly signed measure $\mu,\nu$. The sigma algebra of the product measure space $X\times Y$ is
\[\sigma\qty(\{A\times B\mid A\in\calA\land B\in\calB\}),\]
denoted as $\calA\otimes\calB$, and the possibly signed measure of it is $\mu\times\nu$ restricted to $\calA\otimes\calB$.
\sssc{\text{\textsigma}-finite measure}
Let $X$ be a set, $\mathcal{R}$ be a ring of sets on $X$, and $\mu$ being a signed or positive pre-measure. $\mu$ is \text{\textsigma}-finite iff there exists an index family $(A_n)_{n\in\mathbb{N}}\subseteq\calR$ such that
\[\forall n\in\bbN\colon\mu(A_n)<\infty\land\bigcup_{n\in\mathbb{N}}A_n=X.\]
If $\mu$ is a \text{\textsigma}-finite measure, the measure space $(X,\Sigma,\mu)$ is called a \text{\textsigma}-finite measure space.
\sssc{Measures agree on generating \textpi-system agree theorem}
Suppose $\mu_1,\mu_2$ are two measures on a measurable space $(X,\calF)$, and $\calF$ is generated by a \textpi-system $P$ on $X$. If
\[\forall A\in P\colon\mu_1(A)=\mu_2(A)\]
and
\[\mu_1(X)=\mu_2(X)<\infty,\]
then $\mu_1=\mu_2$.
\bpr
Define $D=\{A\in\calF\colon\mu_1(A)=\mu_2(A)\}$, which is a \textlambda-system since $\mu_1(X)=\mu_2(X)<\infty$. By \textpi-\textlambda theorem, $\sigma(P)\subseteq D$.
\epr
\sssc{Carathéodory's extension theorem}
Let $X$ be a set, $\mathcal{R}$ be a ring of sets on $X$, $\mu_0$ be a pre-measure on $\mathcal{R}$, and $\sigma(\mathcal{R})$ be the \text{\textsigma}-algebra generated by $\mathcal{R}$.

The Carathéodory's extension theorem states that there exists a measure $\mu$ on $\sigma(\mathcal{R})$ such that $\mu_0$ is an extension of $\mu$.

Moreover, if $\mu_0$ is \text{\textsigma}-finite, then the extension $\mu$ is unique and also \text{\textsigma}-finite.
\bpr
PLACEHOLDER
\epr
\sssc{Measurable functions}
Let $(X,\Sigma)$ and $(Y,T)$ be two measurable spaces. A function $f\colon D\subseteq X\to Y$ is said to be measurable iff for every $E\in T$, the preimage
\[f^{-1}(E)\in\Sigma.\]
\sssc{Borel measurable functions}
Let $(X,\calB)$ be a Borel measurable space and $(Y,T)$ be a measurable spaces. A function $f\colon D\subseteq X\to Y$ is said to be Borel measurable iff for every $E\in T$, the preimage
\[f^{-1}(E)\in\calB.\]
\sssc{Almost everywhere (a.e.)}
If $(X,\Sigma ,\mu)$ is a measure space, a property $P$ is said to hold almost everywhere or $\mu$-almost everywhere in $Y\subseteq X$ if there exists a measurable set $N\in \Sigma\land N\subseteq Y$ with $\mu (N)=0$, such that for all $x\in Y\setminus N$, the property $P$ holds at $x$.
\sssc{Continuity from below}
Let $(A_n)_{n\in\bbN}$ be a sequence of measurable sets $A_{i-1}\subseteq A_i$, and let $A=\bigcup_{n=1}^{\infty}A_n$. Then a measure $\mu$ satisfies
\[\mu(A)=\lim_{n\to\infty}\mu(A_n).\]
\bpr
Define
\[B_1=A_1,\]
\[B_n=A_n\setminus A_{n-1},\quad n\ge 2.\]
\[\mu(A_n)=\sum_{k=1}^n\mu(B_k),\]
\[\mu(A)=\sum_{k=1}^{\infty}\mu(B_k).\]
\epr
\sssc{Continuity from above}
Let $(A_n)_{n\in\bbN}$ be a sequence of measurable sets $A_{i-1}\superseteq A_i$, and let $A=\bigcap_{n=1}^{\infty}A_n$. Then a measure $\mu$ such that $\mu(A_1)<\infty$ satisfies
\[\mu(A)=\lim_{n\to\infty}\mu(A_n).\]
\bpr
Define
\[B_n=A_1\setminus A_n.\]
By continuity from below,
\[\mu(\bigcup_{k=1}^nB_n)=\mu(A_1\setminus A)=\lim_{n\to\infty}\mu(A_1\setminus A_n)\]
Since $A,A_n\subseteq A_1$, by countable additivity,
\[\mu(A_1\setminus A)=\mu(A_1)-\mu(A)\]
\[\mu(A_1\setminus A_n)=\mu(A_1)-\mu(A_n)\]
\[\mu(A_1)-\mu(A)=\lim_{n\to\infty}(\mu(A_1)-\mu(A_n))=\mu(A_1)-\lim_{n\to\infty}\mu(A_n)\]
Since $\mu(A_1)<\infty$,
\[\mu(A)=\lim_{n\to\infty}\mu(A_n)\]
\epr
\ssc{Borel measure}
\sssc{Borel measure
A Borel measure is a measure on a Borel \text{\textsigma}-algebra.
\sssc{Radon measure}
Let $X$ be a locally compact Hausdorff space, and $\mu$ be a Borel measure of $X$.
\begin{itemize}
\item $\mu$ is called inner regular or tight iff, for every open set $U$, $\mu(U)$ equals the supremum of $\mu(K)$ over all compact subsets $K$ of $U$.
\item $\mu$ is called outer regular iff, for every Borel set $B$, $\mu(B)$ equals the infimum of $\mu(U)$ over all open sets $U$ that contain $B$.
\item $\mu$ is called locally finite iff every point of $X$ has a neighborhood $U$ for which $\mu(U)$ is finite.
\item $\mu$ is called a Radon measure if it is inner regular and locally finite.
\end{itemize}
\ssc{Lebesgue measure (勒貝格測度)}
\sssc{Real numbers}
Let $R$ be the set of all open intervals, and $\ell((b-a))=b-a$ be the length function. For any subset $E\subseteq\mathbb{R}$, the Lebesgue outer measure $\lambda(E)$ is defined as
\[\lambda(E)=\inf\left\{\sum_{k=1}^{\infty}\ell(I_k)\colon(I_k)_{k\in\mathbb{N}}\subseteq R\land E\subseteq\bigcup_{k=1}^{\infty}I_k\right\}.\]
\sssc{Real vector space}
Let $R$ be the set of all open rectangle in $\bbR^n$ with $n\in\bbN$, that is, $\prod_{i=1}^nI_i$ for open intervals $I_i$. For any $C\in R$, the Lebesgue outer measure $\lambda(C)$ is defined as
\[\lambda(C)=\prod_{i=1}^n\ell(I_i).\]
For any subset $E\subseteq\bbR^n$ with $n\in\bbN$, the Lebesgue outer measure $\lambda(E)$ is defined as
\[\lambda(E)=\inf\left\{\sum_{k=1}^{\infty}\lambda(C_k)\colon(C_k)_{k\in\mathbb{N}}\subseteq R\land E\subseteq\bigcup_{k=1}^{\infty}C_k\right\}.\]
\sssc{Carathéodory criterion}
A set $E\subseteq\mathbb{R}^n$ satisfies the Carathéodory criterion iff
\[\forall A\subseteq\bbR^n\colon\lambda(A)=\lambda(A\cap E)+\lambda(A\cap (\mathbb{R}^n\setminus E)).\]
A set $E$ that satisfies the Carathéodory criterion is Lebesgue-measurable, with its Lebesgue measure being defined as its Lebesgue outer measure. The set of all such $E$ is the Lebesgue \text{\textsigma}-algebra.

$\calB(\bbR^n)$ is a subset of Lebesgue \text{\textsigma}-algebra.

A set $E\subseteq\mathbb{R}^n$ that does not satisfy the Carathéodory criterion is not Lebesgue-measurable. ZFC proves that such sets do exist.
\bpr
PLACEHOLDER
\epr
\sssc{Unique Borel measure that is translation invariant}
Lebesgue measure is the unique Borel measure that is translation invariant, up to scaling by a positive constant.
\bpr
PLACEHOLDER
\epr
\ssc{Hausdorff measure (郝斯多夫測度)}
\sssc{Hausdorff outer measure}
Let $(X,d)$ be a metric space. For any subset $U\subseteq X$, let $\operatorname{diam}(U)$ denote its diameter, that is, $\operatorname{diam}(\varnothing)=0$ and for all $U\ne\varnothing$,
\[\operatorname{diam}(U)=\sup\{d(x,y)\mid x,y\in U\}.\]
Let $S$ be any subset of $X$, and $\delta$ be a positive real number. Define
\[H_{\delta}^d(S)=\inf\left\{\sum_{k=1}^\infty\qty(\operatorname{diam}(U_k))^d\mid S\subseteq\bigcup_{k=1}^\infty U_i\subseteq X\land\forall k\in\bbN\colon\operatorname{diam}(U_k)<\delta\right\}.\]
Note that $H_\delta^d(S)$ is nonincreasing in $\delta$. Thus, 
$\lim_{\delta\to 0}H_\delta^d(S)\in[0,\infty]$. The Hausdorff outer measure $H^d(S)$ is defined as
\[H^d(S):=\lim_{\delta\to 0}H_\delta^d(S).\]
\sssc{Carathéodory criterion}
A set $E\subseteq X$ satisfies the Carathéodory criterion iff
\[\forall A\subseteq X\colon H^d(A)=H^d(A\cap E)+H^d(A\cap (X\setminus E)).\]
A set $E$ that satisfies the Carathéodory criterion is Hausdorff-measurable, with its Hausdorff measure being defined as its Hausdorff outer measure. The set of all such $E$ is the Hausdorff \text{\textsigma}-algebra.

$\calB(X)$ is a subset of Hausdorff \text{\textsigma}-algebra.
\bpr
PLACEHOLDER
\epr

A set $E\subseteq X$ that does not satisfy the Carathéodory criterion is not Hausdorff-measurable.
\ssc{Hahn decomposition theorem}
\sssc{Hahn decomposition theorem}
PLACEHOLDER
\sssc{Jordan decomposition theorem}
PLACEHOLDER
\end{document}